% !TEX program = pdflatex
\documentclass[14pt, russian]{scrartcl}
\let\counterwithout\relax
\let\counterwithin\relax
\usepackage{float}
\usepackage{xcolor}
\usepackage{extsizes}
\usepackage{subfig}
\usepackage[export]{adjustbox}
\usepackage{tocvsec2}
\usepackage[subfigure]{tocloft}
\usepackage[newfloat]{minted}

\makeatletter
\AddToHook{begindocument/before}{\@ifpackageloaded{minted}{\removefromtoclist[float]{lol}}{}}
\makeatother

\AtBeginEnvironment{figure}{\vspace{0.5cm}}
\AtBeginEnvironment{table}{\vspace{0.5cm}}
\AtBeginEnvironment{listing}{\vspace{0.5cm}}
\AtBeginEnvironment{algorithm}{\vspace{0.5cm}}
\AtBeginEnvironment{minted}{\vspace{-0.5cm}}

\usepackage{fancyvrb}
\usepackage{ulem,bm,mathrsfs,ifsym}
\normalem
\renewcommand{\emph}[1]{#1}
\usepackage{sectsty}

%%% Поля и разметка страницы %%%
\usepackage{pdflscape}
\usepackage{geometry}
\geometry{a4paper,tmargin=2cm,bmargin=2cm,lmargin=3cm,rmargin=1cm}

%%% Математические пакеты %%%
\usepackage{amsthm,amsfonts,amsmath,amssymb,amscd}
\usepackage{mathtools}
\usepackage[perpage]{footmisc}

\KOMAoptions{fontsize=14pt}

\makeatletter
\def\showfontsize{\f@size{} point}
\makeatother

\makeatletter
\setlength{\@fptop}{0pt}
\makeatother

\usepackage{tempora}
\usepackage{cmap}
\usepackage[T1]{fontenc}
\usepackage[utf8]{inputenc}
\usepackage[english, main=russian]{babel}
\usepackage{indentfirst}

%%% Таблицы %%%
\usepackage{longtable}
\usepackage{multirow,makecell,array}
\usepackage{booktabs}

%%% Общее форматирование
\usepackage{soulutf8}
\usepackage{icomma}

%%% Изображения %%%
\usepackage{graphicx}
\usepackage{wrapfig}

%%% Списки %%%
\usepackage{enumitem}

%%% Подписи %%%
\usepackage{caption}

\addto\captionsrussian{%
  \renewcommand{\listingname}{Листинг}%
}
%%% Счётчики %%%
\usepackage[figure,table,section]{totalcount}
\DeclareTotalCounter{lstlisting}
\usepackage{totcount}
\usepackage{totpages}

\linespread{1.3}

%%% Гиперссылки %%%
\usepackage{hyperref}

%%% Выравнивание и переносы %%%
\sloppy
\clubpenalty=10000
\widowpenalty=10000

\makeatletter
\let\@@scshape=\scshape
\renewcommand{\scshape}{%
  \ifnum\strcmp{\f@series}{bx}=\z@
    \usefont{T1}{cmr}{bx}{sc}%
  \else
    \ifnum\strcmp{\f@shape}{it}=\z@
      \fontshape{scsl}\selectfont
    \else
      \@@scshape
    \fi
  \fi}
\makeatother

\usepackage{ifthen}
\newcounter{intvl}
\newcounter{otstup}
\newcounter{contnumeq}
\newcounter{contnumfig}
\newcounter{contnumtab}
\newcounter{pgnum}
\newcounter{bibliosel}
\newcounter{chapstyle}
\newcounter{headingdelim}
\newcounter{headingalign}
\newcounter{headingsize}
\newcounter{tabcap}
\newcounter{tablaba}
\newcounter{tabtita}

\setcounter{intvl}{3}
\setcounter{otstup}{0}
\setcounter{contnumeq}{1}
\setcounter{contnumfig}{1}
\setcounter{contnumtab}{1}
\setcounter{pgnum}{0}
\setcounter{bibliosel}{1}
\setcounter{chapstyle}{1}
\setcounter{headingdelim}{1}
\setcounter{headingalign}{0}
\setcounter{headingsize}{0}
\setcounter{tabcap}{0}
\setcounter{tablaba}{2}
\setcounter{tabtita}{1}

\DeclareCaptionLabelSeparator*{emdash}{~--- }
\captionsetup[listing]{position=top,labelsep=colon,justification=centering,singlelinecheck=false}
\captionsetup[lstlisting]{position=top,labelsep=colon,justification=centering,singlelinecheck=false}
\captionsetup[figure]{labelsep=emdash,font=onehalfspacing,position=bottom}

\ifthenelse{\equal{\thetabcap}{0}}{%
    \newcommand{\tabcapalign}{\raggedright}
}

\DeclareCaptionFormat{tablecaption}{\tabcapalign #1#2#3}
\captionsetup[table]{labelsep=emdash}
\captionsetup[table]{format=tablecaption,singlelinecheck=off,font=onehalfspacing,position=top,skip=-5pt}
\DeclareCaptionLabelFormat{continued}{Продолжение таблицы~#2}
\setlength{\belowcaptionskip}{.2cm}
\setlength{\intextsep}{0ex}

\renewcommand{\thesubfigure}{\asbuk{subfigure}}
\captionsetup[subfigure]{font={normalsize},
    labelformat=brace,
    justification=centering,
}

\definecolor{linkcolor}{rgb}{0.0,0,0}
\definecolor{citecolor}{rgb}{0,0.0,0}
\definecolor{urlcolor}{rgb}{0,0,0}

\hypersetup{
    linktocpage=true,
    plainpages=true,
    colorlinks,
    linkcolor={linkcolor},
    citecolor={citecolor},
    urlcolor={urlcolor},
    pdflang={ru},
}
\urlstyle{same}

\setlength{\parindent}{2.5em}

\renewcommand{\labelitemi}{---}
\setlist{nosep,
    labelindent=\parindent,leftmargin=*
}

\usepackage{ragged2e}
\usepackage{needspace}
\usepackage[explicit]{titlesec}
\usepackage{placeins}
\usepackage{xparse}
\usepackage{csquotes}

\usepackage{listingsutf8}
\usepackage{url}
\usepackage{algorithm, algorithmicx}
\usepackage[noend]{algpseudocode}
\usepackage{blkarray}
\usepackage{chngcntr}
\usepackage{tabularx}
\usepackage{tikz}
\usetikzlibrary{shapes.geometric, arrows.meta, positioning}

% --- Библиография ---
\usepackage[backend=biber,
    bibstyle=gost-numeric,
    citestyle=nature]{biblatex}
\addbibresource{biblio.bib}
% Как в РПЗ: [Электронный ресурс] у @online (поле media = {eresource} в .bib)
\DefineBibliographyStrings{russian}{%
  mediaeresource = {{Электронный ресурс}},%
}

\titleformat{name=\section,numberless}[block]{\normalfont\Large\centering}{}{0em}{#1}
\titleformat{\section}[block]{\normalfont\Large\bfseries\raggedright}{}{0em}{\thesection\hspace{0.25em}#1}
\titleformat{\subsection}[block]{\normalfont\Large\bfseries\raggedright}{}{0em}{\thesubsection\hspace{0.25em}#1}
\titleformat{\subsubsection}[block]{\normalfont\large\bfseries\raggedright}{}{0em}{\thesubsubsection\hspace{0.25em}#1}

\let\Algorithm\algorithm
\renewcommand\algorithm[1][]{\Algorithm[#1]\setstretch{1.5}}

\usepackage{pifont}
\usepackage{calc}
\usepackage{suffix}
\DeclareQuoteStyle{russian}
    {\guillemotleft}{\guillemotright}[0.025em]
    {\quotedblbase}{\textquotedblleft}
\ExecuteQuoteOptions{style=russian}
\newcommand{\enq}[1]{\enquote{#1}}
\newcommand{\eng}[1]{\begin{english}#1\end{english}}

\definecolor{mygreen}{rgb}{0,0.6,0}
\definecolor{mygray}{rgb}{0.5,0.5,0.5}
\definecolor{mymauve}{rgb}{0.58,0,0.82}
\renewcommand{\listalgorithmname}{Список алгоритмов}
\floatname{algorithm}{Листинг}
\renewcommand{\lstlistingname}{Листинг}
\renewcommand{\thealgorithm}{\arabic{algorithm}}

\newcommand{\refAlgo}[1]{(листинг \ref{#1})}
\newcommand{\refImage}[1]{(рисунок \ref{#1})}

\renewcommand{\theenumi}{\arabic{enumi}.}
\renewcommand{\labelenumi}{\arabic{enumi}.}
\renewcommand{\theenumii}{\arabic{enumii}}
\renewcommand{\labelenumii}{(\arabic{enumii})}
\renewcommand{\theenumiii}{\roman{enumiii}}
\renewcommand{\labelenumiii}{(\roman{enumiii})}
\renewcommand{\labelitemi}{---}
\renewcommand{\labelitemii}{---}

\lstdefinelanguage{GoLike}{
  keywords={package,const,type,var,func,function,struct,class,dynvar,if,else,while,for,return,int,char,nil,weak,true,false},
  sensitive=true,
  morecomment=[l]{//},
  morecomment=[s]{/*}{*/},
  commentstyle=\color{mygreen},
  morestring=[b]",
  stringstyle=\color{mymauve}
}

\lstdefinelanguage{Refal}{
  sensitive=true,
  morecomment=[s]{/*}{*/},
  commentstyle=\color{mygreen},
  morestring=[b]',
  stringstyle=\color{mymauve}
}

\makeatletter
\def\p@subsection{}
\def\p@subsubsection{\thesection\,\thesubsection\,}
\makeatother
\newcommand{\prog}[1]{{\ttfamily\small#1}}
\lstset{
  backgroundcolor=\color{white},
  basicstyle=\ttfamily\footnotesize,
  breakatwhitespace=false,
  breaklines=true,
  captionpos=top,
  commentstyle=\color{mygreen},
  extendedchars=true,
  inputencoding=utf8,
  frame=single,
  keepspaces=true,
  keywordstyle=\normalfont,
  language=GoLike,
  numbers=left,
  numbersep=5pt,
  xleftmargin=25pt,
  xrightmargin=25pt,
  numberstyle=\small\color{black},
  rulecolor=\color{black},
  showspaces=false,
  showstringspaces=false,
  showtabs=false,
  stepnumber=1,
  stringstyle=\color{mymauve},
  tabsize=8,
  title=\lstname
}

\newcommand{\anonsection}[1]{\cleardoublepage
\phantomsection
\addcontentsline{toc}{section}{\protect\numberline{}#1}
\section*{#1}\vspace*{2.5ex}}
\newcommand{\sectionbreak}{\clearpage}
\renewcommand{\sectionfont}{\normalsize}
\renewcommand{\cftsecleader}{\cftdotfill{\cftdotsep}}

\renewcommand{\cftsecfont}{\normalfont\large}
\renewcommand{\cftsubsecfont}{\normalfont\large}

\setlength{\cftsecindent}{0pt}
\setlength{\cftsubsecindent}{0pt}
\setlength{\cftsubsubsecindent}{0pt}

\usepackage{caption}
\usepackage{amsmath}
\usepackage{amsfonts}
\usepackage{mathtools}
    \DeclarePairedDelimiter\abs{\lvert}{\rvert}
    \DeclarePairedDelimiter\norm{\lVert}{\rVert}
\DeclareTextCommandDefault{\textvisiblespace}{%
  \mbox{\kern.06em\vrule \@height.3ex}%
  \vbox{\hrule \@width.3em}%
  \hbox{\vrule \@height.3ex}}
\newsavebox{\spacebox}
\begin{lrbox}{\spacebox}
\verb*! !
\end{lrbox}
\newcommand{\aspace}{\usebox{\spacebox}}
\DeclareTotalCounter{listing}

\makeatletter
\renewcommand*{\p@subsubsection}{}
\makeatother

\begin{document}
\sloppy

\def\figurename{Рисунок}

\begin{titlepage}
\thispagestyle{empty}
\newpage

\vspace*{-30pt}
\hspace{-45pt}
\begin{minipage}{0.17\textwidth}
\hspace*{-20pt}\centering
\includegraphics[width=1.3\textwidth]{Emblem.png}
\end{minipage}
\begin{minipage}{0.82\textwidth}\small \textbf{
\vspace*{-0.7ex}
\hspace*{-10pt}\centerline{Министерство науки и высшего образования Российской Федерации}
\vspace*{-0.7ex}
\centerline{Федеральное государственное автономное образовательное учреждение }
\vspace*{-0.7ex}
\centerline{высшего образования}
\vspace*{-0.7ex}
\centerline{<<Московский государственный технический университет}
\vspace*{-0.7ex}
\centerline{имени Н.Э. Баумана}
\vspace*{-0.7ex}
\centerline{(национальный исследовательский университет)>>}
\vspace*{-0.7ex}
\centerline{(МГТУ им. Н.Э. Баумана)}}
\end{minipage}

\vspace{-2pt}
\hspace{-34.5pt}\rule{\textwidth}{2.5pt}

\vspace*{-20.3pt}
\hspace{-34.5pt}\rule{\textwidth}{0.4pt}

\vspace{0.5ex}
\noindent \small ФАКУЛЬТЕТ\hspace{80pt} <<Информатика и системы управления>>

\vspace*{-16pt}
\hspace{35pt}\rule{0.855\textwidth}{0.4pt}

\vspace{0.5ex}
\noindent \small КАФЕДРА\hspace{50pt} <<Теоретическая информатика и компьютерные технологии>>

\vspace*{-16pt}
\hspace{25pt}\rule{0.875\textwidth}{0.4pt}

\vspace{3em}

\begin{center}
\Large \bf{РАСЧЕТНО-ПОЯСНИТЕЛЬНАЯ ЗАПИСКА\\\textbf{\textit{К НАУЧНО-ИССЛЕДОВАТЕЛЬСКОЙ РАБОТЕ\\НА ТЕМУ:}} \\}
\end{center}

\vspace*{-6ex}
\begin{center}
\Large
% Тема: как в rpz_go_nuiyap.tex — линии 0.9\textwidth, 1.2pt; после строки текста \vspace*{-3ex}, чтобы базовая линия «примыкала» к правилу; две нижние пустые линии — плотный стек как в РПЗ.
\textit{\textbf{«\,Компилятор базисного подмножества}}\par
\vspace*{-3ex}
\rule{0.9\textwidth}{1.2pt}

\vspace*{0.55ex}
\textit{\textbf{Рефала-5 с представлением данных,}}\par
\vspace*{-3ex}
\rule{0.9\textwidth}{1.2pt}

\vspace*{0.55ex}
\textit{\textbf{использующим раскрученные очереди\,»}}\par
\vspace*{-3ex}
\rule{0.9\textwidth}{1.2pt}

\vspace*{-0.2ex}
\rule{0.9\textwidth}{1.2pt}

\vspace*{-0.2ex}
\rule{0.9\textwidth}{1.2pt}
\end{center}

\vspace{\fill}

\newlength{\ML}
\settowidth{\ML}{«\underline{\hspace{0.7cm}}» \underline{\hspace{2cm}}}

\noindent Студент \underline{\hspace{1.5cm}} \hfill \underline{\hspace{4cm}}\quad
\underline{\hspace{4cm}}

\vspace{-2.1ex}
\noindent\hspace{9ex}\scriptsize{(Группа)}\normalsize\hspace{170pt}\hspace{2ex}\scriptsize{(Подпись, дата)}\normalsize\hspace{30pt}\hspace{6ex}\scriptsize{(И.О. Фамилия)}\normalsize

\bigskip

\noindent Руководитель  \hfill \underline{\hspace{4cm}}\quad
\underline{\hspace{4cm}}

\vspace{-2ex}
\noindent\hspace{13.5ex}\normalsize\hspace{170pt}\hspace{2ex}\scriptsize{(Подпись, дата)}\normalsize\hspace{30pt}\hspace{6ex}\scriptsize{(И.О. Фамилия)}\normalsize
\vfill

\begin{center}
\textsl{\the\year{} г.}
\end{center}
\end{titlepage}

\setlength{\tabcolsep}{3pt}
\newpage
\setcounter{page}{2}


\renewcommand\contentsname{\hfill{\normalfont{СОДЕРЖАНИЕ}}\hfill}
\tableofcontents
\newpage

\anonsection{ВВЕДЕНИЕ}

Язык Рефал-5 относится к семейству языков, основанных на переписывании термов:
программа задаётся набором правил сопоставления образца с объектным выражением и
построения результирующего выражения. На практике стоимость вычислений в таком подходе
часто определяется не только числом применённых правил, но и эффективностью
представления объектных выражений и операций над ними (вставка/удаление с концов,
конкатенация).

Цель НИР: исследовать и обосновать применение раскрученных двусторонних
очередей для представления объектных выражений при реализации базисного подмножества
Рефала-5, а также выполнить проектирование программной системы, реализующей выбранное
подмножество языка.

Для достижения цели решаются задачи:
\begin{itemize}
  \item изучить базовые конструкции Рефала-5 и определить состав базисного подмножества, достаточного для практического использования;
  \item исследовать операции над объектными выражениями, определяющие вычислительную стоимость выполнения правил;
  \item описать раскрученную двустороннюю очередь (мелкий и глубокий уровни представления), её инварианты и принцип ребалансировки;
  \item сравнить предложенное представление с альтернативами и сформулировать область применимости;
  \item спроектировать архитектуру программной системы (входной язык, внутренние представления, модули и интерфейсы).
\end{itemize}

\section{Обзор предметной области}

В разделе кратко описаны базовые конструкции и терминология языка Рефал-5,
необходимые для дальнейшего обсуждения. Основными понятиями являются
объектное выражение, сопоставление с образцом и построение результата правила.
В качестве иллюстраций приведены несколько простых функций, которые будут
использоваться в разделе~\ref{sec:research} при разборе вычислительной стоимости
разных представлений объектных выражений в памяти.

\subsection{Программа, функции и предложения}
\label{sec:refal-program}

Программа на Рефале --- это набор функций. Каждая функция задаётся одним или
несколькими \textbf{предложениями} (правилами) вида \enq{образец = результат},
заключёнными в фигурные скобки. При вызове функции интерпретатор просматривает
её предложения сверху вниз и берёт первое, чей образец сопоставим с
аргументом~\cite{refal_book,refal05_site}. По подстановке, полученной при
сопоставлении, строится результирующее выражение, которое заменяет вызов в
текущем поле зрения.

Простейший пример --- рекурсивная функция, вычисляющая длину объектного
выражения (листинг~\ref{lst:length}):

\begin{lstlisting}[language=Refal,caption={Функция подсчёта длины выражения},label={lst:length}]
Length {
   = 0;
  s.First e.Rest = <Add 1 <Length e.Rest>>;
}
\end{lstlisting}

Функция \prog{Length} имеет два предложения. Первое срабатывает на пустом
выражении и возвращает \prog{0}. Второе сопоставляется с непустым выражением:
оно связывает первый терм с переменной \prog{s.First}, остаток --- с
переменной \prog{e.Rest} и рекурсивно подсчитывает длину остатка.
Имена с префиксами \prog{s.}, \prog{t.}, \prog{e.} обозначают переменные
образца разных типов; их семантика рассматривается ниже.

\subsection{Объектные выражения, термы и символы}

Все данные в Рефале представлены \textbf{объектными выражениями}.
Объектное выражение --- это конечная последовательность \textbf{термов}
(возможно пустая). Каждый терм имеет один из двух видов:

\begin{itemize}
  \item \textbf{атомарный терм} --- собственно \textbf{символ} в
        терминологии Рефала; в Рефале-5 различают три разновидности
        символов: \textbf{символ-слово} (идентификатор: \prog{Found},
        \prog{True}, \prog{Nil}), \textbf{макроцифра} (целое число
        произвольной длины) и \textbf{литера} (один знак строкового
        литерала);
  \item \textbf{скобочный терм} --- объектное выражение, заключённое в
        круглые скобки.
\end{itemize}

Важно не смешивать понятия \enq{символ} и \enq{литера}. Символ-слово ---
это атомарный объект (метка), а не последовательность литер: при
сопоставлении \prog{Found} ведёт себя как один атомарный терм. Запись в
одинарных кавычках, например \prog{'abc'}, обозначает три литеры подряд,
то есть объектное выражение длины три. Аналогично, число \prog{42} --- это
один символ (макроцифра), а не последовательность из двух цифр-литер.

Сопоставление с образцом опирается на три класса \textbf{переменных
образца}:

\begin{itemize}
  \item \(s\)-\textbf{переменная} (\prog{s.X}) связывается с одним
        \emph{символом} --- одним атомарным термом любой из трёх
        разновидностей;
  \item \(t\)-\textbf{переменная} (\prog{t.X}) связывается с одним
        \emph{термом}, в том числе со скобочным;
  \item \(e\)-\textbf{переменная} (\prog{e.X}) связывается с произвольной,
        возможно пустой, цепочкой термов.
\end{itemize}

Простой пример с переменной \(e\)-типа --- функция, возвращающая последний
символ непустой последовательности атомарных термов (листинг~\ref{lst:last}):

\begin{lstlisting}[language=Refal,caption={Функция получения последнего символа},label={lst:last}]
Last {
  e.Init s.Last = s.Last;
}
\end{lstlisting}

При вызове \prog{<Last 1 2 3>} переменная \prog{e.Init} связывается с
выражением \prog{1 2}, а \prog{s.Last} --- с символом \prog{3}. Вариант, в
котором \prog{e.Init} забирала бы всё выражение целиком, не годится:
тогда для \prog{s.Last} не осталось бы ни одного терма, и сопоставление
было бы неудачным. Алгоритм сопоставления учитывает требования к концам
образца и подбирает разбиение, при котором все переменные связываются
согласованно.

\subsection{Неоднозначное сопоставление}
\label{sec:ambiguous-matching}

Если в образце встречается несколько \(e\)-переменных, сопоставление
становится \textbf{неоднозначным}: одно и то же объектное выражение можно
разбить разными способами. Стандартный пример --- функция, проверяющая
вхождение заданного символа в выражение (листинг~\ref{lst:contains}):

\begin{lstlisting}[language=Refal,caption={Неоднозначное сопоставление при поиске символа},label={lst:contains}]
Contains {
  s.Needle e.Left s.Needle e.Right = True;
  s.Needle e.Any                   = False;
}
\end{lstlisting}

Первое предложение требует, чтобы и в начале, и где-то дальше в выражении
встретился один и тот же символ \prog{s.Needle}. Образец содержит две
\(e\)-переменные \prog{e.Left} и \prog{e.Right}, и при сопоставлении
интерпретатор перебирает варианты: он начинает с пустого \prog{e.Left}, и
если очередной \prog{s.Needle} не оказывается на текущей позиции,
удлиняет \prog{e.Left} на один терм и пробует снова. Каждое такое
удлинение --- это комбинация операций \enq{взять один терм с начала
остатка} и \enq{дописать этот терм в конец накапливаемого префикса}, то
есть \emph{отделение начала} и \emph{конкатенация}.

Иными словами, даже логически простое правило приводит к множеству шагов,
на каждом из которых интерпретатор разбивает и собирает фрагменты
выражения. Стоимость отдельного такого шага зависит от того, как
выражение представлено в памяти. Этому вопросу посвящён
раздел~\ref{sec:research}.

\section{Исследование и анализ решений}
\label{sec:research}

В разделе рассматривается, как объектное выражение хранится и преобразуется в
памяти при реализации Рефала-5. Сначала фиксируется набор базовых операций над
выражением (\ref{sec:adt-ops}), который нужен интерпретатору при сопоставлении и
построении результата. Далее обсуждаются распространённые представления объектных
выражений и их практические компромиссы. Завершает раздел описание представления,
предлагаемого в данной работе, и сводное сравнение с альтернативами.

\subsection{Объектное выражение как абстрактный тип}
\label{sec:adt-ops}

Чтобы говорить о стоимости разных представлений, удобно зафиксировать набор
операций, через которые выражается интерпретация любого предложения:

\begin{itemize}
  \item \textbf{проверка на пустоту};
  \item \textbf{отделение начала}: получить первый терм и остаток выражения;
  \item \textbf{отделение конца}: получить последний терм и остаток выражения;
  \item \textbf{конкатенация двух выражений} --- основная операция при
        построении результата правила;
  \item \textbf{выражение из одного терма} (создание сингулярного выражения).
\end{itemize}

Эти операции вызываются на каждом шаге сопоставления и каждом шаге сборки
результата. Например, при переборе \(e\)-переменной
(\ref{sec:ambiguous-matching}) последовательно работают отделение начала
и конкатенация; при возврате результата правила вида \prog{= <F s.X> <G
e.Y>;} итоговое выражение собирается из нескольких фрагментов с помощью
конкатенации. Поэтому выбор представления --- это, по существу, выбор
того, какие из перечисленных операций должны быть дёшевы, а какие могут
стать дорогими.

Ниже сначала рассматриваются два классических списковых представления, в
которых дёшевы операции на концах и конкатенация, а проблемой становится
\emph{копирование} значений переменных. Затем разбираются массивные и
векторно-списковые представления, в которых, наоборот, копирование дёшево, а
\emph{конкатенация} может становиться дорогой. После этого описано
гипотетическое представление на основе деревьев конкатенаций. В конце
предлагается раскрученная двусторонняя очередь, амортизированно
устраняющая обе проблемы.

\subsection{Классическое (плоское двусвязное) представление}
\label{sec:flat-list}

Этот способ представления использовался в большинстве реализаций Рефала
вплоть до 1980-х годов и часто называется \textbf{классическим}~\cite{skorobogatov_repr}
или \textbf{плоским двусвязным}~\cite{abramov_romanenko_1988}. Объектное
выражение хранится как двусвязный список звеньев, где каждое звено
соответствует либо атомарному символу, либо одной из двух структурных
скобок. У звена-скобки есть дополнительное поле, указывающее на парную
скобку. Контейнер выражения хранит указатели на первое и последнее звенья.

Стоимость основных операций при таком представлении:

\begin{itemize}
  \item отделение начала или конца --- \(O(1)\): сдвинуть указатель; если
        крайнее звено --- открывающая скобка, можно прыгнуть через поле
        парной скобки на закрывающую и получить скобочный терм целиком;
  \item конкатенация двух выражений --- \(O(1)\): достаточно перенастроить
        четыре указателя на стыке;
  \item копирование выражения --- \(O(N)\), где \(N\) --- число
        звеньев в записи (включая все вложенные скобки): необходимо
        физически создать новый список звеньев.
\end{itemize}

Дешёвая конкатенация делает классическое представление удобным для
построения результата, однако дорогое копирование порождает классическую
\textbf{проблему копирования}: если переменная встречается в правой части
предложения большее число раз, чем в левой, каждое \enq{лишнее} вхождение
требует полной копии её значения. Стандартная иллюстрация --- функция
подстановки из препринта Абрамова и Романенко~\cite{abramov_romanenko_1988}:
для каждого символа выражения по таблице ищется его замена.

\begin{lstlisting}[language=Refal,caption={Подстановка по таблице (пример проблемы копирования)},label={lst:subst}]
Subst {
  (e.Tab)            = ;
  (e.Tab) s.X e.Rest = <Lookup (e.Tab) s.X> <Subst (e.Tab) e.Rest>;
}
\end{lstlisting}

Переменная \prog{e.Tab} в левой части появляется один раз, а в правой ---
дважды (передаётся в \prog{Lookup} и в рекурсивный вызов \prog{Subst}).
При классическом представлении это означает, что на каждом шаге рекурсии
таблица копируется целиком. Если длина исходного выражения и длина
таблицы оба порядка \(n\), общая стоимость становится \(O(n^2)\), и
попытка программировать в функциональном стиле упирается в эту
квадратичную асимптоту. Распространённый \enq{инженерный} обход --- так
называемый \enq{сквозной просмотр выражения} с двумя стеками, переводящий
\prog{Subst} в императивный по духу вариант без многократного
использования \prog{e.Tab}; такое решение позволяет избежать копирований,
но затемняет код и препятствует распараллеливанию~\cite{abramov_romanenko_1988}.

\subsection{Компромиссное представление с подвешенными скобками}
\label{sec:hanging-list}

Чтобы смягчить проблему копирования, не отказываясь от списковой основы,
было предложено \textbf{компромиссное представление}~\cite{skorobogatov_repr}.
Оно используется в реализации Рефала-6 и ряде учебных проектов, а также
обсуждается в документации Рефала-05~\cite{refal05_site}.

В этом представлении звено списка соответствует не отдельному знаку
записи, а одному \textbf{терму верхнего уровня}. Скобочный терм хранится
особым образом: вместо двух парных звеньев (открывающей и закрывающей
скобок) звено содержит ссылку на отдельный кольцевой список, представляющий
вложенное выражение. У головы кольца хранится \textbf{счётчик ссылок}.

Стоимость операций по сравнению с плоским списком:

\begin{itemize}
  \item отделение начала и конца --- \(O(1)\), как и раньше;
  \item конкатенация --- \(O(1)\);
  \item копирование верхнего уровня --- \(O(N_{\text{верх}})\),
        где \(N_{\text{верх}}\) --- число термов верхнего уровня;
        вложенные выражения при копировании \emph{не дублируются}, у их
        кольцевых списков лишь увеличивается счётчик ссылок.
\end{itemize}

Применительно к функции \prog{Subst} это даёт ощутимый выигрыш в
типичных сценариях: если таблица состоит из пар \prog{(s.Key e.Value)},
то при копировании \prog{e.Tab} физически дублируются только звенья
верхнего уровня, а вложенные пары остаются разделёнными между копиями.

Платой становится более сложное управление памятью: необходимо
поддерживать счётчики ссылок при разрушении и сборке выражений и аккуратно
обрабатывать ситуации, когда переменная связана со \emph{внутренностью}
скобочного терма, который оказался расшарен (тогда модификация одной
\enq{копии} могла бы повредить другие). Кроме того, наличие циклических
связей у голов колец усложняет автоматическую сборку мусора.

\subsection{Массивное представление}
\label{sec:array-repr}

Альтернативный путь, предложенный в работе Абрамова и
Романенко~\cite{abramov_romanenko_1988}, состоит в том, чтобы хранить
объектное выражение в \textbf{непрерывном участке памяти} (массиве
тегированных ячеек). Каждая ячейка содержит либо атомарный символ, либо
открывающую или закрывающую структурную скобку. Скобочный терм при таком
хранении не требует отдельного указателя: он занимает диапазон ячеек от
открывающей до закрывающей скобки.

Главное преимущество массивного представления --- \textbf{дешёвое копирование
переменной}. Значение \(e\)-переменной \prog{e.X} представляется парой
индексов \((l, r)\) внутри массива; копирование переменной --- это просто
копирование пары чисел, то есть \(O(1)\). Это превращает функциональный
стиль программирования в Рефале из \enq{дорогого} в \enq{нормальный}: даже
многократное использование одной и той же переменной в правой части не
приводит к физическому дублированию данных. На функции \prog{Subst} из
\ref{sec:flat-list} массивное представление даёт линейную асимптотику без
каких-либо ухищрений с императивным переписыванием.

Стоимость основных операций:

\begin{itemize}
  \item отделение начала и конца --- \(O(1)\): сдвинуть один из индексов;
  \item копирование переменной --- \(O(1)\);
  \item конкатенация двух выражений длины \(|X|\) и \(|Y|\) ---
        \(O(|X| + |Y|)\): требуется выделить новый массив подходящей
        длины и физически скопировать в него содержимое обоих фрагментов.
\end{itemize}

Так возникает \textbf{проблема конкатенации}, симметричная проблеме
копирования. Если результат правила собирается в цикле, как в типовой
функции отображения

\Needspace{10\baselineskip}
\begin{lstlisting}[language=Refal,caption={Построение результата через рекурсивную конкатенацию},label={lst:map}]
Map {
   = ;
  s.X e.Rest = <F s.X> <Map e.Rest>;
}
\end{lstlisting}

то на \(n\) шагах рекурсии возникают \(n\) конкатенаций. При наивной
массивной реализации каждая из них копирует уже накопленный префикс, и
суммарная стоимость становится \(O(n^2)\).

В препринте~\cite{abramov_romanenko_1988} предлагается смягчение --- так
называемые \textbf{дырки} (запас свободного места) по краям массива. Если
к выражению с дыркой справа дописывается короткий фрагмент длины \(k\),
то вместо создания нового массива достаточно скопировать только эти \(k\)
ячеек в свободную область. На сериях коротких приписываний это даёт
амортизированно линейное общее время. Однако если дырка исчерпывается,
приходится снова выделять новый массив и переносить всё содержимое;
гарантию \(O(1)\) на произвольных последовательностях операций дырки не
дают. Дополнительно в~\cite{abramov_romanenko_1988} обсуждаются меры по
снижению затрат на конкатенацию --- распознавание частных случаев,
функции с несколькими входами и выходами, использование скобок для
структуризации данных.

\subsection{Векторно-списковое представление}
\label{sec:vector-list}

Векторно-списковое представление, предложенное
Скоробогатовым~\cite{skorobogatov_repr}, нацелено на то, чтобы сохранить
дешёвое копирование переменных от массивного представления и одновременно
избежать линейной конкатенации.

Идея состоит в том, чтобы изменить интерфейс между функциями. Аргументы
функций передаются как \textbf{векторы} (массивы тегированных ячеек), но
\textbf{результат} функция возвращает не как один сплошной вектор, а как
\textbf{список ссылок на фрагменты векторов}. Конкатенация фрагментов
откладывается: пока вычисление не дошло до точки, где результат должен
стать пассивным аргументом другой функции (в частности, попасть внутрь
скобочного терма), фрагменты остаются отдельными звеньями.

Применительно к функции \prog{Map} это означает следующее. На каждом шаге
рекурсии вместо физического присоединения \prog{<F s.X>} к результату
\prog{<Map e.Rest>} формируется лёгкое звено \enq{сначала такой-то
фрагмент, потом такой-то}. Когда вычисление функции завершено и её
результат должен стать пассивным выражением, накопленный список
фрагментов собирается в единый вектор за один проход, что даёт линейное
общее время вместо квадратичного.

Платой становится более сложный интерфейс языка сборки и более
ответственное управление временем жизни фрагментов: пока на фрагмент есть
ссылки в незавершённых вычислениях, исходный вектор нельзя освободить или
переиспользовать. Подробное описание представления и связанной системы
команд абстрактной Рефал-машины приведено в~\cite{skorobogatov_repr}.

\subsection{Гипотетическое представление: дерево конкатенаций}
\label{sec:rope-repr}

Среди представлений, не получивших распространения в реализациях Рефала,
интерес представляет \textbf{дерево конкатенаций} (rope). В этом
представлении объектное выражение хранится в виде двоичного дерева, где
листья --- небольшие фрагменты массива, а внутренние узлы соответствуют
операциям конкатенации поддеревьев.

Преимущество rope в том, что конкатенация двух выражений сводится к
созданию одного нового внутреннего узла, что даёт \(O(1)\) с периодической
перебалансировкой за \(O(\log n)\). Копирование переменной также \(O(1)\)
--- это ссылка на корень поддерева. Слабая сторона --- \textbf{операции на
концах}: чтобы получить первый или последний терм, нужно спуститься по
дереву к крайнему листу, что стоит \(O(\log n)\) и зависит от баланса.
Для Рефала это особенно неудобно, потому что \(e\)-переменные требуют как
раз многократного отделения термов с краёв.

В литературе по реализациям Рефала rope обсуждается лишь как теоретическая
возможность~\cite{refal05_site}.

\subsection{Предлагаемое представление: раскрученная двусторонняя очередь}
\label{sec:bootstrapped-deque}

Все рассмотренные выше представления имеют выраженное \enq{слабое место}:
плоский список --- проблему копирования; компромиссный --- неполное её
устранение и сложное управление памятью; массивное --- проблему
конкатенации; векторно-списковое --- сложный интерфейс и привязку к
определённой модели вычисления; rope --- дорогие операции на концах. В
данной работе предлагается использовать \textbf{раскрученную двустороннюю
очередь}, описанную в монографии
Окасаки~\cite{okasaki_pfds}, как структуру, амортизированно решающую обе
ключевые проблемы одновременно.

Идея структуры опирается на два уровня вложенности.

\begin{itemize}
  \item \textbf{Мелкий уровень} (\(Shallow\)) --- буфер фиксированной
        ёмкости \(B\), организованный как циклический массив. Локальные
        операции на концах при наличии места выполняются как сдвиг
        индексов, без выделения памяти.
  \item \textbf{Глубокий уровень} (\(Deep\)) --- тройка \(\langle
        front,\, mid,\, back\rangle\), где крайние компоненты
        (\(front\), \(back\)) --- мелкие буферы, а средний (\(mid\)) ---
        двусторонняя
        очередь мелких буферов. Структура определена рекурсивно: средняя
        очередь сама может быть как мелкой, так и глубокой.
\end{itemize}

Ключевой инвариант: внутренние буферы средней очереди заполнены не менее
чем наполовину. Этот инвариант предотвращает вырождение глубокой
структуры в длинную цепочку почти пустых буферов и обеспечивает, что
глубина рекурсивного представления растёт логарифмически от длины
выражения.

Стоимость основных операций:

\begin{itemize}
  \item отделение начала или конца --- \(O(1)\) амортизированно;
  \item конкатенация двух выражений --- \(O(1)\) амортизированно (за счёт
        перетекания крайних буферов в среднюю очередь);
  \item использование переменной несколько раз --- \(O(1)\), благодаря
        структурному разделению: новая версия выражения
        ссылается на те же неизменяемые буферы и узлы, что и старая.
\end{itemize}

Структурное разделение в данном случае --- инженерная техника, а не
свойство языка: в самой модели Рефала-5 нет изменяемых ячеек памяти, и
вычисление концептуально строит новые выражения, а задача реализации в
том, чтобы за этой моделью не стояло физическое копирование.

\subsubsection{Структура мелкого и глубокого уровней}

Рисунок~\ref{fig:shallow_ring} показывает идею мелкого буфера: элементы
размещены циклически, концы задаются индексами \enq{head} и
\enq{tail}; в общем случае head может находиться правее tail, и
разрыв \enq{обходится} по модулю ёмкости \(B\).

\begin{figure}[!htb]
\centering
\begin{tikzpicture}[
    cell/.style={rectangle, draw, minimum width=0.9cm, minimum height=0.55cm, font=\footnotesize, align=center},
    lab/.style={font=\footnotesize},
]
  \node[cell] (c0) { };
  \node[cell, right=0cm of c0] (c1) { };
  \node[cell, right=0cm of c1] (c2) {x};
  \node[cell, right=0cm of c2] (c3) {y};
  \node[cell, right=0cm of c3] (c4) {z};
  \node[cell, right=0cm of c4] (c5) { };
  \node[cell, right=0cm of c5] (c6) { };
  \node[cell, right=0cm of c6] (c7) { };
  \node[cell, right=0cm of c7] (c8) { };
  \node[lab, below=0.25cm of c2] {head};
  \node[lab, below=0.25cm of c4] {tail};
\end{tikzpicture}
\caption{Идея мелкого буфера: элементы размещены циклически, концы задаются индексами}
\label{fig:shallow_ring}
\end{figure}

Глубокий уровень состоит из двух мелких буферов и средней очереди, которая
сама хранит мелкие буферы (рисунок~\ref{fig:deep_deque}). Рекурсивность
означает, что средняя очередь при необходимости тоже становится глубокой,
и мы получаем структуру с логарифмической по основанию \(B/2\) глубиной.

\begin{figure}[!htb]
\centering
\begin{tikzpicture}[
    box/.style={rectangle, draw, rounded corners, minimum width=3.2cm, minimum height=0.85cm, align=center, font=\small},
    arrow/.style={-{Stealth[length=2.5mm]}, thick}
]
    \node[box] (front) {front\\(shallow buffer)};
    \node[box, right=1.2cm of front] (mid) {mid\\(deque of buffers)};
    \node[box, right=1.2cm of mid] (back) {back\\(shallow buffer)};
    \draw[arrow] (front) -- (mid);
    \draw[arrow] (mid) -- (back);
\end{tikzpicture}
\caption{Глубокий уровень раскрученной двусторонней очереди: два мелких буфера и средняя очередь буферов}
\label{fig:deep_deque}
\end{figure}

\subsubsection{Краткое обоснование амортизированных оценок}

Большинство операций затрагивает только крайние буферы \(front\) и
\(back\): пока в них есть свободное место, добавление и удаление
сводятся к сдвигу индексов в циклическом массиве и являются строго
\(O(1)\). Перетекание буфера в среднюю очередь происходит только при
достижении границ заполненности и требует \(O(B)\). Если каждой такой
дорогой операции приписать стоимость, накопленную на предшествующих
локальных операциях, то средняя стоимость шага остаётся около константы;
это и есть классический амортизированный аргумент~\cite{okasaki_pfds}.

Интуитивно редкие глубокие переносы можно сравнить с цепочкой переносов
при прибавлении единицы в позиционной системе счисления: вероятность
цепочки длины \(k\) убывает геометрически, и вклад в \emph{среднюю}
стоимость образует сходящийся ряд. Для основания \(q \ge 2\) ожидаемое
число переносов
\[
  \mathbb{E}[K] = \sum_{k\ge 1} \frac{k(q-1)}{q^{k+1}} = \frac{1}{q-1}.
\]
Это иллюстрация, а не строгое доказательство для конкретной реализации.
Полные техники амортизированного анализа функциональных структур данных
(включая раскрутку и ленивые вычисления) систематически изложены у
Окасаки~\cite{okasaki_pfds}.

\subsubsection{Сводное сравнение представлений}

Сравнение рассмотренных представлений приведено в
таблице~\ref{tab:repr-comparison}. Зафиксирована стоимость трёх ключевых
операций, через которые выражается основная вычислительная нагрузка
интерпретатора: отделение терма с конца, конкатенация двух выражений и
использование одного значения переменной несколько раз.

\begin{table}[!htb]
\caption{\centering Сравнение представлений объектных выражений по стоимости основных операций}
\label{tab:repr-comparison}
\small
\centering
\begin{tabular}{@{}p{4.2cm}p{2.4cm}p{3.0cm}p{3.4cm}@{}}
\toprule
\textbf{Представление} & \textbf{Концы} & \textbf{Конкатенация} & \textbf{Копирование переменной} \\
\midrule
Классическое (плоское двусвязное) & \(O(1)\) & \(O(1)\) & \(O(N)\) (вся запись) \\
Компромиссное (подвешенные скобки) & \(O(1)\) & \(O(1)\) & \(O(N_{\text{верх}})\) \\
Массивное (с дырками) & \(O(1)\) & \(O(|X|+|Y|)\), линейная амортизация при последовательных приписываниях & \(O(1)\) \\
Векторно-списковое & \(O(1)\) & откладывается, \(O(1)\) на сборке & \(O(1)\) \\
Дерево конкатенаций (rope) & \(O(\log n)\) & \(O(1)\), перебалансировка \(O(\log n)\) & \(O(1)\) \\
Раскрученная двусторонняя очередь & \(O(1)\) аморт. & \(O(1)\) аморт. & \(O(1)\) \\
\bottomrule
\end{tabular}
\end{table}

\subsubsection{Область применимости и ограничения}

Раскрученная двусторонняя очередь оправдана прежде всего тогда, когда
выражения достаточно длинные, а в правилах активно используются
\(e\)-переменные, порождающие многократные разбиения и сборки. На
коротких выражениях и в программах, где почти каждое правило связано с
одним-двумя термами, выигрыш от сложной структуры может уступить простым
спискам или массивам из-за более высокого постоянного множителя.

Целевой реализацией в данной работе является среда JVM. Это естественно
сочетается с неизменяемыми объектами и структурным разделением, но
накладные расходы на работу сборщика мусора нужно учитывать при выборе
ёмкости \(B\) мелкого буфера: слишком малое \(B\) приводит к большому
числу мелких объектов, слишком большое --- ухудшает локальность
модификаций функционального стиля.

\subsection{Выводы по разделу}

Требования к представлению объектных выражений базисного подмножества
Рефала-5 сводятся к следующему: дёшевые операции на концах, дешёвая
конкатенация при построении результата и отсутствие физического
копирования при многократном использовании одной переменной. Существующие
представления закрывают часть этих требований: классическое и
компромиссное списковые --- быстрые на концах и при конкатенации, но
дорогие на копировании; массивное и векторно-списковое --- быстрые на
копировании, но имеют ограничения по конкатенации или по интерфейсу
языка сборки; rope --- быстрая конкатенация и копирование, но дорогие
операции на концах. Раскрученная двусторонняя очередь покрывает все три
требования амортизированно за \(O(1)\) и поэтому выбрана в качестве
представления объектных выражений в проектируемой системе.

\section{Проектирование программной системы}
\label{sec:design}

В третьем разделе фиксируется архитектура программной системы: состав модулей, интерфейсы и данные,
которые они обрабатывают. Проектирование опирается на выбранное представление выражений (раскрученная двусторонняя очередь)
и на грамматику базисного подмножества.

\subsection{Функциональные требования}

К системе предъявляются требования:
\begin{itemize}
  \item приём исходного текста на подмножестве Рефала-5;
  \item выдача сообщений об ошибках с привязкой к позиции во входном тексте;
  \item выполнение программы в модели \enq{выбор первого применимого правила} для функций;
  \item сокращённая библиотека встроенных функций (только необходимый минимум по постановке).
\end{itemize}

\subsection{Архитектура и потоки данных}

Система включает цепочку этапов: подготовка входа (при необходимости), лексический анализ, синтаксический анализ,
семантический анализ и выполнение правил. Представление объектных выражений едино для этапов выполнения и используется
при сопоставлении и сборке результата.

На практике удобно разрешить внешний \enq{рассахариватель}, который переводит исходный текст полного Рефала-5 к синтаксису
базисного подмножества (условия, блоки и др.), чтобы не расширять грамматику минимальной системы. Такой шаг не меняет
семантику рассматриваемых программ при условии корректности рассахаривания, но выносит сложные конструкции из ядра
транслятора.

\begin{figure}[!htb]
\centering
\begin{tikzpicture}[
    node distance=0.55cm and 0.55cm,
    box/.style={rectangle, draw, rounded corners, minimum width=2.35cm, minimum height=0.78cm, align=center, font=\footnotesize},
    data/.style={rectangle, draw, dashed, minimum width=2.15cm, minimum height=0.62cm, align=center, font=\scriptsize},
    arrow/.style={-{Stealth[length=2.2mm]}, thick}
]
    \node[box] (prep) {Подготовка\\входа};
    \node[box, right=of prep] (lexer) {Лексический\\анализ};
    \node[box, right=of lexer] (parser) {Синтаксический\\анализ};
    \node[box, right=of parser] (sem) {Семантический\\анализ};
    \node[box, right=of sem] (exec) {Исполнение\\правил};

    \node[data, below=0.55cm of prep] (srcfull) {Исходный текст};
    \node[data, below=0.55cm of lexer] (srcbasis) {Базисный текст};
    \node[data, below=0.55cm of parser] (toks) {Токены};
    \node[data, below=0.55cm of sem] (ast) {Дерево программы};
    \node[data, below=0.55cm of exec] (expr) {Выражения};

    \draw[arrow] (srcfull) -- (prep);
    \draw[arrow] (prep) -- (srcbasis);
    \draw[arrow] (srcbasis) -- (lexer);
    \draw[arrow] (lexer) -- (toks);
    \draw[arrow] (toks) -- (parser);
    \draw[arrow] (parser) -- (ast);
    \draw[arrow] (ast) -- (sem);
    \draw[arrow] (sem) -- (exec);
    \draw[arrow] (exec) -- (expr);
\end{tikzpicture}
\caption{Этапы обработки программы и основные данные}
\label{fig:arch}
\end{figure}

Целевая реализация рассматривается в экосистеме JVM (например, Java): это естественно сочетается с представлением
программы и выражений как ссылочных структур и со сборкой мусора, но накладные расходы GC должны учитываться при выборе
гранулярности объектов (термы, буферы раскрученной двусторонней очереди, узлы дерева программы).

\subsection{Внутренние представления данных}

В проектировании фиксируются следующие структуры:
\begin{itemize}
  \item \textbf{токены} с координатами (строка/столбец) для точной диагностики;
  \item \textbf{дерево программы} как набор функций и правил сопоставления;
  \item \textbf{объектные выражения} как двусторонняя очередь термов.
\end{itemize}

\subsection{Интерфейс контейнера выражений}

Для независимости модулей удобно выделить интерфейс контейнера выражения, содержащий операции:
проверку пустоты, отделение терма с начала/конца, добавление/удаление на концах, итерацию и конкатенацию.
Конкретная реализация (мелкий и глубокий уровни раскрученной двусторонней очереди) скрыта за этим
интерфейсом и может заменяться без изменения логики сопоставления.

\subsection{Декомпозиция на модули}

С практической точки зрения систему удобно декомпозировать на набор модулей со строго определёнными ответственностями:
\begin{itemize}
  \item \textbf{модуль ввода} --- чтение текста программы, выбор кодировки, единый входной поток символов;
  \item \textbf{подготовка входа} (опционально) --- запуск внешнего рассахаривателя и получение текста базисного подмножества;
  \item \textbf{лексический анализатор} --- преобразование потока символов в поток токенов с координатами;
  \item \textbf{синтаксический анализатор} --- построение дерева программы (функции, предложения, выражения);
  \item \textbf{семантический анализатор} --- проверки, которые удобнее выполнить до исполнения (например, определённость функций и согласованность переменных);
  \item \textbf{контейнер выражений} --- представление объектных выражений и операции над ним;
  \item \textbf{механизм сопоставления} --- сопоставление образца и выражения, вычисление подстановки;
  \item \textbf{исполнитель правил} --- выбор первого применимого предложения и построение результата;
  \item \textbf{таблица встроенных функций} --- набор примитивов и диспетчеризация вызовов.
\end{itemize}

Такая декомпозиция снижает связанность: развитие входного языка (например, уточнение синтаксиса) не затрагивает слой контейнера выражений,
а оптимизация представления выражений не затрагивает логику синтаксического анализа.

\subsection{Модель выполнения и точки расширения}

Модель выполнения фиксируется как \enq{вызов функции \(\rightarrow\) просмотр предложений сверху вниз \(\rightarrow\) первое успешное сопоставление}.
Отдельные точки расширения в проектировании:
\begin{itemize}
  \item добавление новых классов термов (при необходимости расширить подмножество);
  \item расширение механизма сопоставления (например, введение дополнительных ограничений или условий);
  \item расширение библиотеки встроенных функций без изменения ядра сопоставления;
  \item возможное добавление оптимизаций (индексация предложений по первым термам, мемоизация вызовов).
\end{itemize}

Даже если оптимизации не входят в круг основных задач данного этапа, их перечисление как точек расширения полезно, поскольку показывает,
что выбранный интерфейс контейнера выражений не препятствует дальнейшему развитию системы.

\subsection{Нефункциональные требования}

Помимо функциональности, для системы важны следующие свойства:
\begin{itemize}
  \item \textbf{предсказуемость потребления памяти}: представление выражений должно использовать структурное разделение и избегать полного копирования;
  \item \textbf{диагностика}: сообщения об ошибках должны указывать координаты во входном тексте и причину ошибки;
  \item \textbf{расширяемость}: состав встроенных функций допускается сокращать/расширять без изменения ядра сопоставления;
  \item \textbf{переносимость}: предпочтителен язык реализации со сборкой мусора и поддержкой обобщённых типов
        (в том числе варианты на JVM).
\end{itemize}

\subsection{Обработка ошибок и форматы сообщений}

Ошибки целесообразно разделять на классы: лексические (неожиданный символ, незакрытый строковый литерал), синтаксические (неожиданная конструкция,
нарушение структуры определения функции) и ошибки выполнения (отсутствие применимого правила). Для повышения качества диагностики
в проектировании закладывается единый формат сообщения: координата (строка, столбец), тип ошибки и краткое описание.

\subsection{Выбор и сокращение встроенной библиотеки}

Сокращённая библиотека встроенных функций проектируется как отдельный слой среды выполнения: каждая встроенная функция имеет строго определённый
интерфейс, а таблица примитивов отделена от механизма сопоставления. Такой дизайн позволяет исключать функции, не нужные для рассматриваемого
подмножества, не затрагивая ядро выполнения.

\subsection{Представление термов и вложенность}

В выбранном подмножестве термы образуют минимально необходимый набор для выражения большинства \enq{структурных} преобразований:
атомарные термы (символы-слова, числа, строковые литералы) и скобочные термы. Скобочный терм содержит вложенное выражение, что фактически вводит
деревоподобную структуру данных поверх линейной последовательности. С точки зрения проектирования это означает, что контейнер выражения
должен хранить не только атомы, но и ссылки на вложенные выражения, причём вложенность может быть произвольной.

Здесь важно различать две роли контейнера:
\begin{itemize}
  \item на \textbf{верхнем уровне} контейнер обслуживает операции над последовательностью термов (концы, конкатенация, разбиение);
  \item на \textbf{уровне скобок} контейнер используется рекурсивно, так как содержимое скобок является самостоятельным выражением.
\end{itemize}
Таким образом, выбранная структура данных применяется единообразно на всех уровнях вложенности, что упрощает архитектуру и повышает
предсказуемость поведения системы.

\subsection{Окружение сопоставления и представление подстановок}

Результатом сопоставления образца и выражения является окружение (подстановка), связывающая переменные образца со значениями.
В проектировании выделяются три класса связываний:
\begin{itemize}
  \item \(s\)-переменная связывается с атомарным термом (число, символ-слово или строковый литерал);
  \item \(t\)-переменная связывается с термом (включая скобочный);
  \item \(e\)-переменная связывается с выражением (цепочкой термов).
\end{itemize}

С практической точки зрения удобно хранить подстановку в виде отображения \enq{имя переменной \(\rightarrow\) значение}, где значение
для \(e\)-переменных является контейнером выражения. Это позволяет строить результат без распаковки цепочек в списки/массивы:
результат собирается из фрагментов через операции конкатенации и добавления на концах, опираясь на интерфейс контейнера.

\subsection{Контракт механизма сопоставления}

Чтобы разделить \enq{логику сопоставления} и \enq{операции над данными}, проектирование фиксирует контракт модуля сопоставления:
\begin{itemize}
  \item вход: образец (последовательность термов с переменными) и выражение-аргумент;
  \item выход: либо \enq{неуспех}, либо подстановка, корректная относительно ограничений типов переменных;
  \item свойство: если сопоставление успешно, то подстановка совместима со всеми повторными вхождениями переменных (одно имя --- одно значение).
\end{itemize}

Такой контракт позволяет в дальнейшем заменять стратегию перебора границ \(e\)-переменных (например, на более эффективную),
не меняя остальную архитектуру. При этом контейнер выражения должен обеспечивать операции, которые нужны сопоставлению: проверку равенства
фрагментов, получение префикса/суффикса и сборку новых выражений с разделением структуры.

\subsection{Сценарии использования системы}

Проектирование исходит из основного сценария --- \textbf{пакетный запуск}: чтение программы из файла, запуск входной функции
с заданным выражением и вывод результата. Такой сценарий опирается на единый формат представления выражений и единый механизм диагностики
и соответствует типичному применению инструмента как отдельной программы.

\textbf{Интерактивный режим} (последовательные запуски с разными входными выражениями без повторной загрузки программы) в рамках
данного проектирования не фиксируется; его можно рассматривать как возможную доработку или расширение --- например, для отладки
и исследования поведения правил.

\anonsection{ЗАКЛЮЧЕНИЕ}

В НИР сформулирована мотивация использования раскрученных двусторонних очередей для представления
объектных выражений в реализации базисного подмножества Рефала-5. Описана структура shallow/deep,
инварианты заполненности и причины, по которым достигаются хорошие амортизированные оценки на концах.
Выполнено проектирование программной системы и исследовательское обоснование выбора структуры данных
для представления объектных выражений.

Проведённое исследование показывает, что выбор контейнера для объектных выражений нельзя рассматривать как второстепенную деталь:
в языке переписывания стоимость вычисления во многом определяется тем, насколько эффективно выполняются разбиения и сборки выражений.
Подход с раскруткой даёт возможность совместить эффективное версионирование выражений через структурное разделение и хорошие оценки операций на концах,
а также создаёт основу для эффективной конкатенации в сценариях интенсивного построения результата.

Сформулированные в работе требования и проектные решения задают основу для дальнейшей доработки реализации: отдельной задачей
останется по мере необходимости расширять поддерживаемый синтаксис и библиотеку встроенных функций в соответствии
с уточняемыми требованиями к системе и согласованными тестовыми сценариями.

\clearpage
\printbibliography[heading=bibintoc,title={СПИСОК ИСПОЛЬЗОВАННЫХ ИСТОЧНИКОВ}]

\end{document}

